% Options for packages loaded elsewhere
% Options for packages loaded elsewhere
\PassOptionsToPackage{unicode}{hyperref}
\PassOptionsToPackage{hyphens}{url}
\PassOptionsToPackage{dvipsnames,svgnames,x11names}{xcolor}
%
\documentclass[
  letterpaper,
  DIV=11,
  numbers=noendperiod]{scrreprt}
\usepackage{xcolor}
\usepackage{amsmath,amssymb}
\setcounter{secnumdepth}{5}
\usepackage{iftex}
\ifPDFTeX
  \usepackage[T1]{fontenc}
  \usepackage[utf8]{inputenc}
  \usepackage{textcomp} % provide euro and other symbols
\else % if luatex or xetex
  \usepackage{unicode-math} % this also loads fontspec
  \defaultfontfeatures{Scale=MatchLowercase}
  \defaultfontfeatures[\rmfamily]{Ligatures=TeX,Scale=1}
\fi
\usepackage{lmodern}
\ifPDFTeX\else
  % xetex/luatex font selection
  \setmainfont[]{Noto Serif CJK KR}
  \setmathfont[]{Latin Modern Math}
\fi
% Use upquote if available, for straight quotes in verbatim environments
\IfFileExists{upquote.sty}{\usepackage{upquote}}{}
\IfFileExists{microtype.sty}{% use microtype if available
  \usepackage[]{microtype}
  \UseMicrotypeSet[protrusion]{basicmath} % disable protrusion for tt fonts
}{}
\makeatletter
\@ifundefined{KOMAClassName}{% if non-KOMA class
  \IfFileExists{parskip.sty}{%
    \usepackage{parskip}
  }{% else
    \setlength{\parindent}{0pt}
    \setlength{\parskip}{6pt plus 2pt minus 1pt}}
}{% if KOMA class
  \KOMAoptions{parskip=half}}
\makeatother
% Make \paragraph and \subparagraph free-standing
\makeatletter
\ifx\paragraph\undefined\else
  \let\oldparagraph\paragraph
  \renewcommand{\paragraph}{
    \@ifstar
      \xxxParagraphStar
      \xxxParagraphNoStar
  }
  \newcommand{\xxxParagraphStar}[1]{\oldparagraph*{#1}\mbox{}}
  \newcommand{\xxxParagraphNoStar}[1]{\oldparagraph{#1}\mbox{}}
\fi
\ifx\subparagraph\undefined\else
  \let\oldsubparagraph\subparagraph
  \renewcommand{\subparagraph}{
    \@ifstar
      \xxxSubParagraphStar
      \xxxSubParagraphNoStar
  }
  \newcommand{\xxxSubParagraphStar}[1]{\oldsubparagraph*{#1}\mbox{}}
  \newcommand{\xxxSubParagraphNoStar}[1]{\oldsubparagraph{#1}\mbox{}}
\fi
\makeatother


\usepackage{longtable,booktabs,array}
\newcounter{none} % for unnumbered tables
\usepackage{calc} % for calculating minipage widths
% Correct order of tables after \paragraph or \subparagraph
\usepackage{etoolbox}
\makeatletter
\patchcmd\longtable{\par}{\if@noskipsec\mbox{}\fi\par}{}{}
\makeatother
% Allow footnotes in longtable head/foot
\IfFileExists{footnotehyper.sty}{\usepackage{footnotehyper}}{\usepackage{footnote}}
\makesavenoteenv{longtable}
\usepackage{graphicx}
\makeatletter
\newsavebox\pandoc@box
\newcommand*\pandocbounded[1]{% scales image to fit in text height/width
  \sbox\pandoc@box{#1}%
  \Gscale@div\@tempa{\textheight}{\dimexpr\ht\pandoc@box+\dp\pandoc@box\relax}%
  \Gscale@div\@tempb{\linewidth}{\wd\pandoc@box}%
  \ifdim\@tempb\p@<\@tempa\p@\let\@tempa\@tempb\fi% select the smaller of both
  \ifdim\@tempa\p@<\p@\scalebox{\@tempa}{\usebox\pandoc@box}%
  \else\usebox{\pandoc@box}%
  \fi%
}
% Set default figure placement to htbp
\def\fps@figure{htbp}
\makeatother





\setlength{\emergencystretch}{3em} % prevent overfull lines

\providecommand{\tightlist}{%
  \setlength{\itemsep}{0pt}\setlength{\parskip}{0pt}}



 


\usepackage{kotex}
\KOMAoption{captions}{tableheading}
\makeatletter
\@ifpackageloaded{tcolorbox}{}{\usepackage[skins,breakable]{tcolorbox}}
\@ifpackageloaded{fontawesome5}{}{\usepackage{fontawesome5}}
\definecolor{quarto-callout-color}{HTML}{909090}
\definecolor{quarto-callout-note-color}{HTML}{0758E5}
\definecolor{quarto-callout-important-color}{HTML}{CC1914}
\definecolor{quarto-callout-warning-color}{HTML}{EB9113}
\definecolor{quarto-callout-tip-color}{HTML}{00A047}
\definecolor{quarto-callout-caution-color}{HTML}{FC5300}
\definecolor{quarto-callout-color-frame}{HTML}{acacac}
\definecolor{quarto-callout-note-color-frame}{HTML}{4582ec}
\definecolor{quarto-callout-important-color-frame}{HTML}{d9534f}
\definecolor{quarto-callout-warning-color-frame}{HTML}{f0ad4e}
\definecolor{quarto-callout-tip-color-frame}{HTML}{02b875}
\definecolor{quarto-callout-caution-color-frame}{HTML}{fd7e14}
\makeatother
\makeatletter
\@ifpackageloaded{bookmark}{}{\usepackage{bookmark}}
\makeatother
\makeatletter
\@ifpackageloaded{caption}{}{\usepackage{caption}}
\AtBeginDocument{%
\ifdefined\contentsname
  \renewcommand*\contentsname{Table of contents}
\else
  \newcommand\contentsname{Table of contents}
\fi
\ifdefined\listfigurename
  \renewcommand*\listfigurename{List of Figures}
\else
  \newcommand\listfigurename{List of Figures}
\fi
\ifdefined\listtablename
  \renewcommand*\listtablename{List of Tables}
\else
  \newcommand\listtablename{List of Tables}
\fi
\ifdefined\figurename
  \renewcommand*\figurename{Figure}
\else
  \newcommand\figurename{Figure}
\fi
\ifdefined\tablename
  \renewcommand*\tablename{Table}
\else
  \newcommand\tablename{Table}
\fi
}
\@ifpackageloaded{float}{}{\usepackage{float}}
\floatstyle{ruled}
\@ifundefined{c@chapter}{\newfloat{codelisting}{h}{lop}}{\newfloat{codelisting}{h}{lop}[chapter]}
\floatname{codelisting}{Listing}
\newcommand*\listoflistings{\listof{codelisting}{List of Listings}}
\makeatother
\makeatletter
\makeatother
\makeatletter
\@ifpackageloaded{caption}{}{\usepackage{caption}}
\@ifpackageloaded{subcaption}{}{\usepackage{subcaption}}
\makeatother
\usepackage{bookmark}
\IfFileExists{xurl.sty}{\usepackage{xurl}}{} % add URL line breaks if available
\urlstyle{same}
\hypersetup{
  pdftitle={통계학 산책},
  pdfauthor={김재광},
  colorlinks=true,
  linkcolor={blue},
  filecolor={Maroon},
  citecolor={Blue},
  urlcolor={Blue},
  pdfcreator={LaTeX via pandoc}}


\title{통계학 산책}
\author{김재광}
\date{}
\begin{document}
\maketitle

\renewcommand*\contentsname{Table of contents}
{
\setcounter{tocdepth}{2}
\tableofcontents
}

\bookmarksetup{startatroot}

\chapter*{머리말}\label{uxba38uxb9acuxb9d0}
\addcontentsline{toc}{chapter}{머리말}

\markboth{머리말}{머리말}

이 책은 유튜브 채널 \textbf{통계학 산책}의 강의 내용을 바탕으로 엮은
것입니다.

통계학을 전공하지 않은 독자도 읽을 수 있도록, 수식보다 직관적 이해를
우선하되 핵심 논리는 엄밀하게 서술하였습니다. 각 강의는 실제 데이터와
사례에서 출발하여, 그 뒤에 숨어 있는 통계학적 원리를 자연스럽게 드러내는
방식으로 구성하였습니다.

\section*{이 책의 구성}\label{uxc774-uxcc45uxc758-uxad6cuxc131}
\addcontentsline{toc}{section}{이 책의 구성}

\markright{이 책의 구성}

{\def\LTcaptype{none} % do not increment counter
\begin{longtable}[]{@{}ll@{}}
\toprule\noalign{}
강의 & 주제 \\
\midrule\noalign{}
\endhead
\bottomrule\noalign{}
\endlastfoot
제1강 & 출구조사의 추억 --- MAR 가정과 무응답 편향 \\
제2강 & 집회 참가자수 추정 --- 페르미 추정과 현장 샘플링 \\
\end{longtable}
}

강의가 추가될 때마다 이 목록은 업데이트됩니다.

\section*{대상 독자}\label{uxb300uxc0c1-uxb3c5uxc790}
\addcontentsline{toc}{section}{대상 독자}

\markright{대상 독자}

\begin{itemize}
\tightlist
\item
  통계학에 관심 있는 일반 독자
\item
  데이터 분석을 공부하는 학부생
\item
  설문·여론조사 실무자
\end{itemize}

\section*{유튜브 채널}\label{uxc720uxd29cuxbe0c-uxcc44uxb110}
\addcontentsline{toc}{section}{유튜브 채널}

\markright{유튜브 채널}

각 챕터에 대응하는 강의 영상은 아래 재생목록에서 볼 수 있습니다. 책과
영상을 함께 활용하시면 더 효과적입니다.

\href{https://youtube.com/playlist?list=PLpX4d5n0gRfQe0RpCR-jH8b4o0iHPeXaa&si=4MVFil9Cb7rMbr6n}{▶
통계학 산책 재생목록 바로가기}

\bookmarksetup{startatroot}

\chapter{제1강: 출구조사의
추억}\label{uxc81c1uxac15-uxcd9cuxad6cuxc870uxc0acuxc758-uxcd94uxc5b5}

\begin{quote}
``통계학은 데이터에서 진실을 끌어내는 학문이다. 그러나 진실을
끌어내려면, 먼저 데이터가 우리에게 거짓말을 하고 있다는 사실을 알아야
한다.''
\end{quote}

\begin{center}\rule{0.5\linewidth}{0.5pt}\end{center}

\section{출구조사란
무엇인가}\label{uxcd9cuxad6cuxc870uxc0acuxb780-uxbb34uxc5c7uxc778uxac00}

선거 당일, 투표소 앞에는 방송사 조사원들이 서 있다. 투표를 마치고 나오는
유권자들에게 ``어느 당에 투표하셨나요?''라고 묻는다. 이것이
\textbf{출구조사(exit poll)}다.

출구조사의 목적은 하나다. 개표가 시작되기 전에 --- 때로는 투표가 끝나는
순간 --- 당선자를 예측하는 것이다. 방송사들이 밤새 생중계를 이어가면서
``○○ 당선 유력''이라는 자막을 띄울 수 있는 것은 이 조사 덕분이다.

그런데 출구조사는 때로 크게 빗나간다. 왜 그럴까? 단순히 표본이
부족해서일까, 아니면 더 구조적인 이유가 있을까?

이 강의는 그 질문에서 시작한다.

\begin{center}\rule{0.5\linewidth}{0.5pt}\end{center}

\section{2012년 강동갑
사례}\label{uxb144-uxac15uxb3d9uxac11-uxc0acuxb840}

\subsection{데이터}\label{uxb370uxc774uxd130}

2012년 제19대 국회의원 선거, 서울 강동갑 선거구의 출구조사 자료를
살펴보자.

\textbf{표 1.1: 2012 국회의원 선거 강동갑 출구조사 표본 자료}

{\def\LTcaptype{none} % do not increment counter
\begin{longtable}[]{@{}
  >{\raggedright\arraybackslash}p{(\linewidth - 12\tabcolsep) * \real{0.1017}}
  >{\raggedright\arraybackslash}p{(\linewidth - 12\tabcolsep) * \real{0.1356}}
  >{\centering\arraybackslash}p{(\linewidth - 12\tabcolsep) * \real{0.1695}}
  >{\centering\arraybackslash}p{(\linewidth - 12\tabcolsep) * \real{0.2034}}
  >{\centering\arraybackslash}p{(\linewidth - 12\tabcolsep) * \real{0.1017}}
  >{\centering\arraybackslash}p{(\linewidth - 12\tabcolsep) * \real{0.1864}}
  >{\centering\arraybackslash}p{(\linewidth - 12\tabcolsep) * \real{0.1017}}@{}}
\toprule\noalign{}
\begin{minipage}[b]{\linewidth}\raggedright
성별
\end{minipage} & \begin{minipage}[b]{\linewidth}\raggedright
연령대
\end{minipage} & \begin{minipage}[b]{\linewidth}\centering
새누리당
\end{minipage} & \begin{minipage}[b]{\linewidth}\centering
통합민주당
\end{minipage} & \begin{minipage}[b]{\linewidth}\centering
기타
\end{minipage} & \begin{minipage}[b]{\linewidth}\centering
응답 거부
\end{minipage} & \begin{minipage}[b]{\linewidth}\centering
총합
\end{minipage} \\
\midrule\noalign{}
\endhead
\bottomrule\noalign{}
\endlastfoot
남성 & 20-29 & 93 & 115 & 4 & 28 & 240 \\
남성 & 30-39 & 104 & 233 & 8 & 82 & 427 \\
남성 & 40-49 & 146 & 295 & 5 & 49 & 495 \\
남성 & 50- & 560 & 350 & 3 & 174 & 1,087 \\
여성 & 20-29 & 106 & 159 & 8 & 62 & 335 \\
여성 & 30-39 & 129 & 242 & 5 & 70 & 446 \\
여성 & 40-49 & 170 & 262 & 5 & 69 & 506 \\
여성 & 50- & 501 & 218 & 7 & 211 & 937 \\
\textbf{총합} & & \textbf{1,809} & \textbf{1,874} & \textbf{45} &
\textbf{745} & \textbf{4,473} \\
\textbf{개표 결과} & & \textbf{62,489} & \textbf{57,909} &
\textbf{1,624} & & \textbf{122,022} \\
\end{longtable}
}

출구조사에서 샘플로 뽑힌 사람은 총 4,473명이 선정되었다. 하지만 이 중
745명, 약 \textbf{17\%}가 응답을 거부했다. 다른 조사와는 달리 출구조사는
참값을 알게 된다. 개표를 통해 최종 결과가 나오기 때문이다. 따라서 이
표의 맨 아래 줄의 개표 결과와 비교하는 것이 이 분석의 핵심이다.

\subsection{단순 분석의
실패}\label{uxb2e8uxc21c-uxbd84uxc11duxc758-uxc2e4uxd328}

출구조사 자료를 분석하여 각 정당별 최종 득표수를 예측하고자 한다. 이
경우 가장 단순한 방법은 응답 거부자 745명을 제외하고, 응답자
3,728명만으로 지지율을 계산하는 것이다.

\textbf{표 1.2: 단순 분석 결과 (\%)}

{\def\LTcaptype{none} % do not increment counter
\begin{longtable}[]{@{}lccc@{}}
\toprule\noalign{}
방법론 & 새누리당 & 통합민주당 & 기타 \\
\midrule\noalign{}
\endhead
\bottomrule\noalign{}
\endlastfoot
무보정 & 48.5 & 50.3 & 1.2 \\
참값 & 51.2 & 47.5 & 1.3 \\
\end{longtable}
}

결과는 충격적이다. 출구조사는 통합민주당 50.3\%, 새누리당 48.5\%로
예측했지만, 실제 개표 결과는 새누리당 51.2\%, 통합민주당 47.5\%였다.
\textbf{당선자가 뒤바뀐 것이다.}

\begin{center}\rule{0.5\linewidth}{0.5pt}\end{center}

\section{인구학적 보정의 시도와
한계}\label{uxc778uxad6cuxd559uxc801-uxbcf4uxc815uxc758-uxc2dcuxb3c4uxc640-uxd55cuxacc4}

\subsection{성·연령대
보정이란}\label{uxc131uxc5f0uxb839uxb300-uxbcf4uxc815uxc774uxb780}

응답 거부자들을 단순히 제외하는 것이 문제라면, 이들의 투표 성향을
\textbf{예측}해서 채워 넣으면 어떨까? 이것이 인구학적 보정의 기본
아이디어다.

``같은 성별, 같은 연령대의 응답자들과 비슷하게 투표했을 것이다''라고
가정하고, 거부자 수에 해당 그룹의 응답자 지지율을 곱해 예측값을 만드는
것이다.

\textbf{예시: 여성 50대 이상 그룹}

이 그룹의 응답자는 726명(= 501 + 218 + 7)이었고, 응답 거부자는
211명이었다.

\[P(\text{새누리당} \mid \text{여성, 50대 이상}) = \frac{501}{726}\]
\[P(\text{통합민주당} \mid \text{여성, 50대 이상}) = \frac{218}{726}\]
\[P(\text{기타} \mid \text{여성, 50대 이상}) = \frac{7}{726}\]

이 비율을 거부자 211명에 적용하면:

\begin{itemize}
\tightlist
\item
  새누리당 지지자 예측: \(211 \times \frac{501}{726} \approx 145.6\)명
\item
  통합민주당 지지자 예측: \(211 \times \frac{218}{726} \approx 63.4\)명
\item
  기타 예측: \(211 \times \frac{7}{726} \approx 2.0\)명
\end{itemize}

이렇게 채워 넣은 값을 통계학에서는 \textbf{대체값(imputation
value)}이라고 부른다.

\subsection{보정의 결과}\label{uxbcf4uxc815uxc758-uxacb0uxacfc}

\textbf{표 1.3: 인구학적 보정 결과 (\%)}

{\def\LTcaptype{none} % do not increment counter
\begin{longtable}[]{@{}lccc@{}}
\toprule\noalign{}
방법론 & 새누리당 & 통합민주당 & 기타 \\
\midrule\noalign{}
\endhead
\bottomrule\noalign{}
\endlastfoot
무보정 & 48.5 & 50.3 & 1.2 \\
성·연령대 보정 & 49.0 & 49.8 & 1.2 \\
\textbf{참값} & \textbf{51.2} & \textbf{47.5} & \textbf{1.3} \\
\end{longtable}
}

보정 후 새누리당 지지율이 48.5\%에서 49.0\%로 올라갔다. 조금 나아지긴
했지만, \textbf{여전히 당선자 예측은 틀렸다.}

\begin{tcolorbox}[enhanced jigsaw, arc=.35mm, bottomrule=.15mm, bottomtitle=1mm, breakable, colback=white, colbacktitle=quarto-callout-important-color!10!white, colframe=quarto-callout-important-color-frame, coltitle=black, left=2mm, leftrule=.75mm, opacityback=0, opacitybacktitle=0.6, rightrule=.15mm, title=\textcolor{quarto-callout-important-color}{\faExclamation}\hspace{0.5em}{결론}, titlerule=0mm, toprule=.15mm, toptitle=1mm]

단순한 인구학적 변수를 이용한 보정 방법으로는 출구조사 응답 거부자의
편향을 제대로 제거할 수 없다.

\end{tcolorbox}

\begin{center}\rule{0.5\linewidth}{0.5pt}\end{center}

\section{통계학적 접근}\label{uxd1b5uxacc4uxd559uxc801-uxc811uxadfc}

\subsection{기호 정의}\label{uxae30uxd638-uxc815uxc758}

\begin{itemize}
\tightlist
\item
  \(Y\): 지지 정당 (우리가 알고 싶은 변수)
\item
  \(X\): 인구학적 변수 (성별, 연령대 --- 관측 가능)
\item
  \(R\): 응답 협조 지시 변수
\end{itemize}

\[R = \begin{cases} 1 & \text{if 응답} \\ 0 & \text{if 거절} \end{cases}\]

\subsection{MAR 가정과 그
한계}\label{mar-uxac00uxc815uxacfc-uxadf8-uxd55cuxacc4}

성·연령대 보정이 암묵적으로 가정하는 것이 있다. 통계학에서
\textbf{MAR(Missing at Random, 임의 결측)}이라고 부르는 가정이다.

\[P(Y \mid X, R=0) = P(Y \mid X, R=1) \tag{MAR}\]

한국말로 풀면: ``\(X\)가 같다면, 응답 거부자와 응답자의 투표 성향이
같다.''

이 가정이 성립하면 인구학적 보정은 완벽하게 작동한다. 문제는 \textbf{이
가정이 현실에서 성립하지 않을 수 있다}는 것이다.

강동갑 사례에서는 인구학적 보정으로는 충분하지 않았다. 그 이유는 아마도
다음과 같이 정리될수 있을 것이다: 같은 성별, 같은 연령대이더라도
\textbf{새누리당 지지자들이 응답을 더 많이 거부했다.} 그 결과 응답자 중
새누리당 지지자의 비율이 실제보다 낮게 잡혔고, 집계 결과가 왜곡되었다.

\begin{center}\rule{0.5\linewidth}{0.5pt}\end{center}

\section{응답 거부자 예측
모형}\label{uxc751uxb2f5-uxac70uxbd80uxc790-uxc608uxce21-uxbaa8uxd615}

\subsection{베이즈 정리의
활용}\label{uxbca0uxc774uxc988-uxc815uxb9acuxc758-uxd65cuxc6a9}

베이즈 정리를 이용하면 \(P(Y \mid X, R=0)\)과 \(P(Y \mid X, R=1)\)을
각각 쓸 수 있다.

\[P(Y \mid X, R=0) = \frac{P(R=0 \mid X,Y)\,P(Y \mid X)}{P(R=0 \mid X)} \tag{1}\]

\[P(Y \mid X, R=1) = \frac{P(R=1 \mid X,Y)\,P(Y \mid X)}{P(R=1 \mid X)} \tag{2}\]

식 (1)을 식 (2)로 나누면 \(P(Y \mid X)\)가 약분되어 다음을 얻는다.

\[\frac{P(Y \mid X, R=0)}{P(Y \mid X, R=1)} = \underbrace{\frac{P(R=0 \mid X,Y)}{P(R=1 \mid X,Y)}}_{\displaystyle\text{Nonresponse odds}} \cdot \underbrace{\frac{P(R=1 \mid X)}{P(R=0 \mid X)}}_{\displaystyle Y\text{와 무관}} \tag{3}\]

\textbf{두 번째 항은 \(Y\)와 무관하다.} 따라서 응답 거부자와 응답자의
투표 성향 차이를 결정하는 것은 오직 첫 번째 항, \textbf{Nonresponse
odds}다.

\subsection{예측 공식의
도출}\label{uxc608uxce21-uxacf5uxc2dduxc758-uxb3c4uxcd9c}

Nonresponse odds 함수를 다음과 같이 정의하자.

\[O(X, Y) \equiv \frac{P(R=0 \mid X,Y)}{P(R=1 \mid X,Y)}\]

그러면 식 (3)으로부터 다음의 \textbf{예측 공식}을 유도할 수 있다.

\[P(Y=y \mid x, R=0) = P(Y=y \mid x, R=1) \cdot \frac{O(x,y)}{E[O(x,Y) \mid x, R=1]} \tag{4}\]

응답 거부자의 투표 성향은, 응답자의 투표 성향을 출발점으로 삼아
\textbf{Nonresponse odds 비율만큼 재조정}한 것이다.

\begin{center}\rule{0.5\linewidth}{0.5pt}\end{center}

\section{Exponential Tilting Model}\label{exponential-tilting-model}

\subsection{로지스틱 모형
가정}\label{uxb85cuxc9c0uxc2a4uxd2f1-uxbaa8uxd615-uxac00uxc815}

응답 협조 확률이 \textbf{로지스틱 회귀 모형}을 따른다고 가정하자.

\[P(R=1 \mid X, Y) = \frac{\exp(\beta_0 + \beta_1 X + \beta_2 Y)}{1 + \exp(\beta_0 + \beta_1 X + \beta_2 Y)}\]

이것을 식 (4)에 대입하면 \(\beta_0\)과 \(\beta_1\) 항이 상쇄되어 다음을
얻는다.

\[P(Y=y \mid x, R=0) = P(Y=y \mid x, R=1) \cdot \frac{\exp(-\beta_2 y)}{E[\exp(-\beta_2 Y) \mid x, R=1]} \tag{5}\]

핵심은 \textbf{\(\beta_2\) 하나}다. 이 모형은 Kim and Yu (2011)에서 처음
소개되었으며, \textbf{Exponential tilting model}이라고 부른다.

\subsection{β₂의 해석}\label{ux3b2ux2082uxc758-uxd574uxc11d}

{\def\LTcaptype{none} % do not increment counter
\begin{longtable}[]{@{}clc@{}}
\toprule\noalign{}
\(\beta_2\)의 값 & 의미 & 편향 방향 \\
\midrule\noalign{}
\endhead
\bottomrule\noalign{}
\endlastfoot
\(= 0\) & MAR 성립, 인구학적 보정으로 충분 & 없음 \\
\(> 0\) & 새누리당 지지자가 응답을 더 잘 함 & 새누리당 과대 추정 \\
\(< 0\) & 새누리당 지지자가 응답을 더 많이 거부 & 새누리당 과소 추정 \\
\end{longtable}
}

\begin{tcolorbox}[enhanced jigsaw, arc=.35mm, bottomrule=.15mm, bottomtitle=1mm, breakable, colback=white, colbacktitle=quarto-callout-note-color!10!white, colframe=quarto-callout-note-color-frame, coltitle=black, left=2mm, leftrule=.75mm, opacityback=0, opacitybacktitle=0.6, rightrule=.15mm, title=\textcolor{quarto-callout-note-color}{\faInfo}\hspace{0.5em}{Note}, titlerule=0mm, toprule=.15mm, toptitle=1mm]

2012년 강동갑의 상황은 \(\beta_2 < 0\)의 경우였다. 따라서 식 (5)에서
새누리당 지지 확률은 첫번째 항 \(P(Y=y | X, R=1)\)에 1보다 더 큰 값을
곱해주는 결과를 가져오게 된다.

\end{tcolorbox}

\begin{center}\rule{0.5\linewidth}{0.5pt}\end{center}

\section{실증 분석 결과}\label{uxc2e4uxc99d-uxbd84uxc11d-uxacb0uxacfc}

\subsection{강동갑 선거구}\label{uxac15uxb3d9uxac11-uxc120uxac70uxad6c}

Riddles, Kim, and Im (2016)은 식 (5)의 \(\beta_2\)를 \textbf{EM
알고리즘}으로 추정하는 방법론을 개발하고 강동갑 데이터에 적용했다.

\textbf{표 1.4: 강동갑 출구조사 분석 결과 (\%)}

{\def\LTcaptype{none} % do not increment counter
\begin{longtable}[]{@{}lccc@{}}
\toprule\noalign{}
방법론 & 새누리당 & 통합민주당 & 기타 \\
\midrule\noalign{}
\endhead
\bottomrule\noalign{}
\endlastfoot
무보정 & 48.5 & 50.3 & 1.2 \\
성·연령대 보정 & 49.0 & 49.8 & 1.2 \\
\textbf{제안된 방법론} & \textbf{51.0} & \textbf{47.7} & \textbf{1.2} \\
참값 & 51.2 & 47.5 & 1.3 \\
\end{longtable}
}

제안된 방법론은 참값 51.2\%에 \textbf{오차 0.2\%포인트} 이내로 접근했다.

\subsection{서울 전체 48개
지역구}\label{uxc11cuxc6b8-uxc804uxccb4-48uxac1c-uxc9c0uxc5eduxad6c}

\textbf{표 1.5: 서울 전체 출구조사 분석 결과 (48개 지역구, 1위 예측
횟수)}

{\def\LTcaptype{none} % do not increment counter
\begin{longtable}[]{@{}lccc@{}}
\toprule\noalign{}
방법론 & 새누리당 & 통합민주당 & 기타 \\
\midrule\noalign{}
\endhead
\bottomrule\noalign{}
\endlastfoot
무보정 & 10 & 36 & 2 \\
성·연령대 보정 & 10 & 36 & 2 \\
\textbf{제안된 방법론} & \textbf{15} & \textbf{29} & \textbf{4} \\
참값 & 16 & 30 & 2 \\
\end{longtable}
}

무보정과 인구학적 보정은 결과가 동일했다. 두 방법 모두 MAR 가정에
의존하기 때문이다. 제안된 방법론은 참값(새누리당 16석)에 훨씬 가까운
15석을 예측했다.

\begin{center}\rule{0.5\linewidth}{0.5pt}\end{center}

\section{마무리}\label{uxb9c8uxbb34uxb9ac}

이 연구의 출발점은 저자의 2006년 출구조사 자문단 참여였다. 당시는 연세대
응용통계학과에 근무할 시절이라 한국의 기관들과의 협력이 수월했다. SBS
방송사의 연락으로 자문단에 참여하게 되었는데 이번 무응답 문제에 대해
알게 되었다. MAR 가정이 성립하지 않는다는 것을 알고 이를 어떻게 해결할수
있을지를 고민하기 시작했다.

그후 미국으로 이민을 오게 되었고 당시의 문제를 좀더 깊이있게 고민할
기회가 생겼다. 그 연구 결과는 Kim and Yu (2011)과 Wang et al.~(2014)으로
이어졌다. 두 논문은 저자의 논문 중 가장 많이 인용되는 논문이 되었고,
MAR이 성립하지 않는 어려운 상황에서도 합리적인 추론이 가능하다는
돌파구를 제시했다는 점에서 의미가 있다.

마지막으로 한 가지 생각을 공유하고 싶다.

출구조사의 성공을 ``1등을 맞추었느냐''로 판단하는 것은 방송사의
관점이다. 통계학자의 관점은 다르다.

\begin{tcolorbox}[enhanced jigsaw, arc=.35mm, bottomrule=.15mm, bottomtitle=1mm, breakable, colback=white, colbacktitle=quarto-callout-tip-color!10!white, colframe=quarto-callout-tip-color-frame, coltitle=black, left=2mm, leftrule=.75mm, opacityback=0, opacitybacktitle=0.6, rightrule=.15mm, title=\textcolor{quarto-callout-tip-color}{\faLightbulb}\hspace{0.5em}{통계학자의 관점}, titlerule=0mm, toprule=.15mm, toptitle=1mm]

\textbf{1등을 틀리더라도 오차가 1\textasciitilde2\%인 예측}이,
\textbf{1등을 맞추더라도 오차가 10\% 이상인 예측}보다 통계학적으로 훨씬
바람직하다.

단순히 결과를 맞추는 것보다, \textbf{결과를 얻어내는 과정이
합리적인가}를 더 중요하게 봐야 한다.

\end{tcolorbox}

\begin{center}\rule{0.5\linewidth}{0.5pt}\end{center}

\section{참고문헌}\label{uxcc38uxace0uxbb38uxd5cc}

\protect\phantomsection\label{refs}

\bookmarksetup{startatroot}

\chapter{제2강: 집회 참가자수
추정}\label{uxc81c2uxac15-uxc9d1uxd68c-uxcc38uxac00uxc790uxc218-uxcd94uxc815}

\begin{quote}
``데이터는 답을 주지 않는다. 올바른 질문을 던질 때 비로소 답이 보인다.''
\end{quote}

\begin{center}\rule{0.5\linewidth}{0.5pt}\end{center}

\section{같은 집회, 다른
숫자}\label{uxac19uxc740-uxc9d1uxd68c-uxb2e4uxb978-uxc22buxc790}

2016년 11월 12일, 광화문 일대에서 대규모 촛불집회가 열렸다. 그날
연합뉴스의 헤드라인은 이렇게 전했다.

\begin{quote}
``사상 최대 촛불집회\ldots{} 주최측-경찰 추산인원 차이 여전''
\end{quote}

주최 측은 \textbf{100만 명}, 경찰은 \textbf{26만 명}. 같은 집회, 같은
날, 같은 장소를 두고 네 배에 가까운 차이가 났다. 왜 그럴까? 그리고 어느
쪽이 더 맞을까?

이 강의는 그 질문에서 시작한다.

\begin{center}\rule{0.5\linewidth}{0.5pt}\end{center}

\section{두 가지 추산
방식}\label{uxb450-uxac00uxc9c0-uxcd94uxc0b0-uxbc29uxc2dd}

\subsection{주최 측: 면적 기반 누적
인원}\label{uxc8fcuxcd5c-uxce21-uxba74uxc801-uxae30uxbc18-uxb204uxc801-uxc778uxc6d0}

주최 측의 방식은 두 단계로 이루어진다.

첫째, 집회 장소의 전체 면적(\(m^2\))을 구한다. 둘째, 단위 면적당 밀집
인원을 곱한다. 여기에 시간대별 유동 인구를 누적하여 집회 전체 시간 동안
현장을 거쳐 간 총인원을 발표한다.

\subsection{경찰 측: 스냅샷 기반 실시간
집계}\label{uxacbduxcc30-uxce21-uxc2a4uxb0c5uxc0f7-uxae30uxbc18-uxc2e4uxc2dcuxac04-uxc9d1uxacc4}

경찰 측은 집회가 가장 밀집된 \textbf{특정 시점}을 기준으로, 경찰이 통제
가능한 주요 거점 위주로 인원을 집계한다.

\subsection{측정 대상 자체가
다르다}\label{uxce21uxc815-uxb300uxc0c1-uxc790uxccb4uxac00-uxb2e4uxb974uxb2e4}

두 방식의 근본적인 차이는 \textbf{무엇을 세느냐}에 있다. 주최 측은 집회
시간 전체를 거쳐 간 \textbf{누적 인원}을, 경찰 측은 특정 시점에 한 곳에
있는 \textbf{동시 인원}을 측정한다. 두 수치가 다른 것은 당연하다.

\begin{center}\rule{0.5\linewidth}{0.5pt}\end{center}

\section{페르미 추정법과 그
한계}\label{uxd398uxb974uxbbf8-uxcd94uxc815uxbc95uxacfc-uxadf8-uxd55cuxacc4}

면적 기반 집계를 \textbf{페르미 추정법}이라 부른다. 공식은 간단하다.

\[\text{전체 인원} = \text{전체 면적} \times \text{단위 면적당 밀도}\]

문제는 \textbf{밀도 설정}에 있다. 주최 측은 1\(m^2\)당
3\textasciitilde4명으로, 경찰 측은 1\textasciitilde2명으로 가정했다.
가정 하나의 차이가 최종 추정치를 몇 배나 벌려놓는다.

페르미 추정은 빠르고 직관적이지만, 핵심 가정이 주관적 판단에 의존한다는
근본적인 한계를 가진다.

\begin{center}\rule{0.5\linewidth}{0.5pt}\end{center}

\section{지하철 승객 통계
이용법}\label{uxc9c0uxd558uxcca0-uxc2b9uxac1d-uxd1b5uxacc4-uxc774uxc6a9uxbc95}

페르미 추정의 한계를 보완하려는 시도 중 하나가 지하철 승객 통계를
활용하는 방법이다. 2016년 11월 13일 연합뉴스는 다음과 같이 보도했다.

\begin{quote}
광화문광장과 서울광장 인근 지하철역 12곳의 집회 당일 하차 승객은
\textbf{86만 4,596명}으로, 전년도 같은 요일 평균보다 \textbf{51만
4,940명} 많았다. 지하철 수송분담률(약 39\%)을 고려하면 집회 참석자는 총
\textbf{132만 명}으로 추산된다.
\end{quote}

공식으로 쓰면 이렇다.

\[514{,}940 \div 0.39 \approx 1{,}320{,}359\]

\subsection{논리 구조를
해부해보면}\label{uxb17cuxb9ac-uxad6cuxc870uxb97c-uxd574uxbd80uxd574uxbcf4uxba74}

이 추론에는 두 단계가 있다.

\begin{enumerate}
\def\labelenumi{\arabic{enumi}.}
\tightlist
\item
  전체 하차 승객 86만 명 → 집회 참석 지하철 이용자 51만 5천 명 (전년
  대비 증가분)
\item
  집회 참석 지하철 이용자 51만 5천 명 → 전체 집회 참석자 132만 명
  (수송분담률 적용)
\end{enumerate}

두 번째 단계에서 핵심 가정이 등장한다. \textbf{집회 참석자의 39\%가
지하철을 이용했다}는 것이다. 그런데 이 39\%는 2014년의 서울시 전체 평균
수송분담률이다. 대규모 집회가 열리는 날, 많은 사람이 대중교통을
선택한다. 이 날의 지하철 이용률이 평상시와 같다는 가정은 현실적으로
적절하지 않다.

\begin{tcolorbox}[enhanced jigsaw, arc=.35mm, bottomrule=.15mm, bottomtitle=1mm, breakable, colback=white, colbacktitle=quarto-callout-important-color!10!white, colframe=quarto-callout-important-color-frame, coltitle=black, left=2mm, leftrule=.75mm, opacityback=0, opacitybacktitle=0.6, rightrule=.15mm, title=\textcolor{quarto-callout-important-color}{\faExclamation}\hspace{0.5em}{핵심 문제}, titlerule=0mm, toprule=.15mm, toptitle=1mm]

평상시의 지하철 수송분담률 39\%를 집회 당일에 그대로 적용하는 것은
타당하지 않다. \textbf{이 비율 자체를 데이터로부터 추정해야 한다.}

\end{tcolorbox}

\begin{center}\rule{0.5\linewidth}{0.5pt}\end{center}

\section{통계학적 접근: 현장
샘플링}\label{uxd1b5uxacc4uxd559uxc801-uxc811uxadfc-uxd604uxc7a5-uxc0d8uxd50cuxb9c1}

\subsection{아이디어}\label{uxc544uxc774uxb514uxc5b4}

해결책은 명쾌하다. 집회 당일의 지하철 이용 비율을 \textbf{현장에서 직접
조사}하여 추정 공식에 반영하자는 것이다.

\subsection{Plan A vs Plan B}\label{plan-a-vs-plan-b}

1,000명을 조사할 예산 하에서 두 가지 플랜이 검토되었다.

\textbf{Plan A}: 조사원 전원을 집회 현장에 배치하여 참석자에게 이동
수단을 묻는다. 모수는 하나: \(P(X=1 \mid Y=1)\) (집회 참석자 중 지하철
이용 비율).

\textbf{Plan B}: 조사원을 두 집단으로 나눈다. 한 팀은 지하철 역에서
하차객에게 집회 참석 여부를 묻고(\(P(Y=1 \mid X=1)\) 추정), 다른 팀은
집회 현장에서 참석자에게 이동 수단을 묻는다(\(P(X=1 \mid Y=1)\) 추정).
두 모수를 동시에 추정한다.

최종적으로 \textbf{Plan B}가 채택되었다. 두 모수를 동시에 추정함으로써
이론적으로 더 효율적인 추정이 가능하기 때문이다.

\subsection{추정 공식}\label{uxcd94uxc815-uxacf5uxc2dd}

모집단을 집회 참석 여부 \(Y\)와 지하철 이용 여부 \(X\)로 분류하면 다음
표를 얻는다.

{\def\LTcaptype{none} % do not increment counter
\begin{longtable}[]{@{}lccc@{}}
\toprule\noalign{}
집회 참석 & 지하철 이용 \((X=1)\) & 다른 교통수단 \((X=0)\) & 총합 \\
\midrule\noalign{}
\endhead
\bottomrule\noalign{}
\endlastfoot
참석 \((Y=1)\) & \(N_{11}\) & \(N_{01}\) & \(\textbf{N}_{+1}\) \\
불참석 \((Y=0)\) & \(N_{10}\) & \(N_{00}\) & \(N_{+0}\) \\
총합 & \(\textbf{N}_{1+}\) & \(N_{0+}\) & \(N_{++}\) \\
\end{longtable}
}

우리의 목표는 \(N_{+1}\), 즉 집회 총 참석자 수를 추정하는 것이다. 다음
등식을 이용한다.

\[N_{+1} = N_{++} \cdot P(Y=1) = N_{++} \cdot P(X=1) \cdot \frac{P(Y=1 \mid X=1)}{P(X=1 \mid Y=1)}\]

\(N_{++} \cdot P(X=1) = N_{1+}\) (지하철 총 하차 승객수)는 서울시 공식
통계로 알 수 있다. 따라서:

\[\boxed{N_{+1} = N_{1+} \cdot \frac{P(Y=1 \mid X=1)}{P(X=1 \mid Y=1)}}\]

두 비율 \(P(Y=1 \mid X=1)\)과 \(P(X=1 \mid Y=1)\)을 현장 조사로 추정하면
된다.

\begin{center}\rule{0.5\linewidth}{0.5pt}\end{center}

\section{분석 결과}\label{uxbd84uxc11d-uxacb0uxacfc}

{\def\LTcaptype{none} % do not increment counter
\begin{longtable}[]{@{}lccc@{}}
\toprule\noalign{}
집회 참석 & 지하철 이용 & 다른 교통수단 & 총합 \\
\midrule\noalign{}
\endhead
\bottomrule\noalign{}
\endlastfoot
참석 & 55만 & 23만 & \textbf{78만} \\
불참석 & 35만 & --- & --- \\
총합 & 90만 & --- & --- \\
\end{longtable}
}

Plan B는 ratio 추정의 형태이므로 \textbf{신뢰구간 계산도 가능}하다는
장점이 있다.

흥미로운 뒷이야기가 있다. 이 조사를 의뢰한 언론사는 비용을 지불하고도
\textbf{여론을 의식하여 최종 결과를 공개하지 않았다}. 통계 결과의 투명한
공개가 얼마나 중요한지를 역설적으로 보여주는 사례다.

\subsection{과소추정
가능성}\label{uxacfcuxc18cuxcd94uxc815-uxac00uxb2a5uxc131}

약 78만 명이라는 추정치는 실제보다 다소 \textbf{과소추정}되었을 가능성이
있다. Plan B에서 지하철 역의 혼잡한 환경은 조사 원칙대로 인터뷰를
진행하기 어렵게 만들었고, 이로 인한 \textbf{비표본 오차(non-sampling
error)}가 결과에 영향을 미쳤을 것으로 판단된다.

이는 Plan A와 Plan B의 트레이드오프를 다시 생각하게 한다. Plan A는
조사원을 한 곳에만 배치하므로 조사원 관리가 용이하고 비표본 오차를 줄일
수 있다. \textbf{이론적으로 우수한 설계가 현장에서 반드시 최선의 결과를
보장하지는 않는다.}

\begin{center}\rule{0.5\linewidth}{0.5pt}\end{center}

\section{모의실험: 적정 표본 수
결정}\label{uxbaa8uxc758uxc2e4uxd5d8-uxc801uxc815-uxd45cuxbcf8-uxc218-uxacb0uxc815}

방법론의 정확성을 검증하기 위해 모의실험을 실시했다. 모집단을 설정한 뒤
반복 샘플링을 통해 집회 참석자 수 \(\hat{\theta}\)를 추정하고, 그 오차를
RMSE로 측정했다.

\[\text{RMSE}(\hat{\theta}) = \left[\mathbb{E}\left\{(\hat{\theta} - \theta)^2\right\}\right]^{1/2}\]

결과는 명확했다. RMSE는 표본 크기가 커질수록 단조 감소하며, 총 표본 크기
\textbf{약 1,000명} 수준에서 RMSE가 5\% 임계선 아래로 내려온다.
1,000명이라는 조사 예산이 실용적으로 충분한 정확도를 확보하기 위한
적절한 수치임이 확인된 것이다.

\begin{center}\rule{0.5\linewidth}{0.5pt}\end{center}

\section{마무리: 통계적 접근의
뼈대}\label{uxb9c8uxbb34uxb9ac-uxd1b5uxacc4uxc801-uxc811uxadfcuxc758-uxbf08uxb300}

페르미 추정은 거친 근사치를 제공한다. 그러나 현대 사회에서는 지하철
통계처럼 다양한 \textbf{보조 정보}를 활용할 수 있다. 통계학은 이 보조
정보를 주관적 판단이 아닌 \textbf{데이터로부터 추정된 모형}을 통해
calibration한다. 이것이 핵심이다.

통계적 접근의 세 가지 뼈대:

\begin{enumerate}
\def\labelenumi{\arabic{enumi}.}
\tightlist
\item
  유용한 보조정보
\item
  통계적 모형
\item
  모형의 모수를 \textbf{데이터}로부터 추정
\end{enumerate}

\begin{tcolorbox}[enhanced jigsaw, arc=.35mm, bottomrule=.15mm, bottomtitle=1mm, breakable, colback=white, colbacktitle=quarto-callout-tip-color!10!white, colframe=quarto-callout-tip-color-frame, coltitle=black, left=2mm, leftrule=.75mm, opacityback=0, opacitybacktitle=0.6, rightrule=.15mm, title=\textcolor{quarto-callout-tip-color}{\faLightbulb}\hspace{0.5em}{통계학자의 관점}, titlerule=0mm, toprule=.15mm, toptitle=1mm]

통계적 방법은 설명 가능하고 투명하다. 물론 비용이 수반된다 --- 데이터
수집 비용과 전문가 활용 비용 모두. 결국 \textbf{통계 품질에 대한 의지}가
그 비용을 지불하는 결정으로 이어진다.

그리고 그 의지는 정치적 중립성에 대한 의지와 긴밀히 연결된다. 통계는
사회적 신뢰를 기반으로 하는 일종의 공공재다.

\end{tcolorbox}




\end{document}
